\section{Related Work}
\label{sec:related}

% Structure: (1) the MentorQA foundation we extend; (2) caregiver needs;
% (3) curated non-AI interventions; (4) evaluations of ungrounded LLMs;
% (5) grounded / domain-specific systems; closing positioning paragraph
% + comparison table (Table~\ref{tab:positioning}).

\paragraph{Mentorship-oriented QA over long-form content.}
Our work directly extends \textsc{MentorQA} \citep{bhalerao2026mentorqa},
which introduced mentorship-focused question answering as a distinct task:
answers that provide guidance, reflection, and learning value rather than
isolated facts. \textsc{MentorQA} released 8,990 QA pairs derived from 180
hours of long-form mentorship videos across four languages, and compared
Single-Agent, Dual-Agent, RAG, and Multi-Agent generation pipelines under
controlled conditions, finding that agentic coordination consistently
produces higher-quality mentorship responses, with the largest gains on
complex topics and lower-resource languages. Notably, health proved to be
its most difficult territory: \emph{Physical Health} was the lowest-scoring
topic (attributed to clinically dense language) and the health categories
were the smallest in the dataset, while general podcast-style mentorship
content offers no safety guarantees for health advice. Its authors
explicitly invite extension of the open-source framework to new domains
with expert oversight. MedicalQA answers this call: we retain the
methodological backbone---transcript chunking, multi-agent QA generation,
and mentorship-oriented evaluation---but ground it in vetted expert
geriatric lectures, target dementia caregivers as a concrete, vulnerable
audience, and adopt the trustworthiness requirements that the healthcare
setting demands.

\paragraph{Dementia caregivers' unmet informational needs.}
NLP research on dementia has concentrated on detection and diagnosis;
a recent systematic review finds caregiver support to be one of the least
explored directions in the field \citep{nlpdementia2025}, even as the
dementia-care community identifies LLMs as a promising avenue for caregiver
assistance \citep{llmdementia2024}. The need itself is well documented.
Analysis of over 60,000 posts on the ALZConnected forum shows caregivers
seeking advice on symptoms, medications, and coping in anonymous,
unverified peer spaces \citep{forumanalysis2024}. Interview studies report
caregivers improvising support from generic tools such as ChatGPT, Reddit,
and online forums, while explicitly asking for \emph{credible information
with transparent sourcing} that adapts to their stage in the caregiving
journey \citep{balancing2025}. A recent taxonomy triangulating 52 studies
with caregiver interviews pinpoints the misalignment: existing
informational resources deliver symptom facts rather than the
skill-building and tailored guidance caregivers need, and current AI
chatbots offer shallow, untailored support \citep{taxonomy2026}. Remote
monitoring technologies, though frequently justified by caregiver burden,
likewise leave caregivers' informational needs unaddressed
\citep{remotemonitoring2025}. These studies identify requirements but
build no system that meets them.

\paragraph{Curated caregiver interventions.}
The accepted non-AI intervention model confirms that expert-authored,
digitally delivered education helps caregivers---but does not scale.
The WHO's iSupport programme packages expert caregiving knowledge into a
five-module self-help curriculum adapted across many countries
\citep{isupport}; yet it is static: caregivers cannot ask their own
questions, and each cultural adaptation requires manual expert effort.
A meta-analysis of computerized cognitive training quantifies the
underlying scalability dilemma: professionally supervised digital training
is effective (verbal-memory SMD 0.72) while unsupervised self-administered
versions lose most of the effect (0.21), leading the authors to call
explicitly for incorporating AI into self-administered interventions
\citep{cct2024}. Expert knowledge helps, but only when a professional is
in the loop---motivating systems that deliver expert-level guidance
interactively and at scale.

\paragraph{LLMs as caregiving advisors.}
Pre-LLM conversational support for dementia was scarce and shallow: a
systematic review found only 6 of 505 candidate chatbots relevant, none
grounding their health content in evidence \citep{chatbots2021}. The LLM
era has not closed the reliability gap. When out-of-the-box models were
evaluated against professional gold standards for home-care instruction,
even GPT-4o omitted relevant details in 60--80\% of scenarios and matched
professional-level specificity in only 10\% of responses
\citep{homecare2025}. Caregivers probing a plain GPT-4o chatbot valued
on-demand access but distrusted its uncited, generic answers; the study's
foremost design recommendation is trustworthy information with transparent
source citation \citep{carey2025}. Injecting curated domain knowledge
helps---prompt engineering over 90 hand-written QA pairs significantly
improved actionability, relevance, and satisfaction
\citep{empowering2026,evaluatingllm2025}---but the authors themselves flag
manual prompt curation as unscalable. Beyond content quality, the
NurValues benchmark shows that even frontier LLMs trail human nurses in
aligning with professional care values under adversarial dialogue, and
that medical fine-tuning alone does not fix this \citep{nurvalues2025}.
The verdict across these evaluations is consistent: caregiver-facing LLM
answers must be grounded in vetted expert knowledge, and the grounding
must scale beyond hand-curated prompts.

\paragraph{Grounded and domain-specific care systems.}
The systems closest to ours ground generation in curated care knowledge.
CaLM combines RAG with fine-tuning over caregiving web resources, showing
a 7B domain model can outperform GPT-3.5 for ADRD caregiving answers
\citep{calm2024}. RISE couples RAG over the evidence-based REACH Caregiver
Notebook with a social robot, reaching 87\% backend correctness in a small
user study \citep{rise2026}. In the professional-care domain, NurseLLM
demonstrates the value of specialization---and its multi-agent variant
adds a further +3.62\% accuracy over the single model, independently
validating agentic designs for care QA \citep{nursellm2025}. Yet every
knowledge source in this line of work is short-form written material: web
pages, a structured notebook, synthetic exam questions, or hand-written QA
pairs. None taps long-form expert instruction such as lecture video; none
uses multi-agent coordination for question generation and answer
synthesis; none frames the interaction as mentorship; and none pairs
scalable LLM-as-judge evaluation with human annotation.

\paragraph{Positioning.}
Table~\ref{tab:positioning} situates MedicalQA. Against the grounded
caregiving systems, we contribute \textsc{MentorQA}'s multi-agent
generation and mentorship-oriented evaluation; against \textsc{MentorQA},
we contribute a vetted expert knowledge source (long-form geriatric
lectures), a healthcare domain and audience, and the safety bar that
domain imposes.

\begin{table*}[t]
\centering
\small
\setlength{\tabcolsep}{4pt}
\resizebox{\textwidth}{!}{%
\begin{tabular}{@{}llllll@{}}
\toprule
 & \textbf{MentorQA} & \textbf{CaLM} & \textbf{RISE} & \textbf{NurseLLM} & \textbf{MedicalQA (ours)} \\
\midrule
Knowledge source & Mentorship podcasts & Caregiving web pages & REACH notebook & Synthetic MCQs & \textbf{Expert geriatric lectures} \\
Domain & General mentorship & ADRD caregiving & ADRD caregiving & Nursing exams & \textbf{Dementia caregiving} \\
Grounding & RAG / segment & RAG + fine-tune & RAG & Fine-tune only & \textbf{RAG over lecture chunks} \\
Generation & Multi-agent & Single model & Single model & MA answer selection & \textbf{Multi-agent QA gen.\ + synthesis} \\
Framing & Mentorship & Q\&A & Education & Exam prep & \textbf{Mentorship-style education} \\
Evaluation & LLM judges + human & Auto metrics & Small user study & MCQ accuracy & \textbf{LLM-judge + human pilot} \\
\bottomrule
\end{tabular}}
\caption{MedicalQA compared with its foundation (\textsc{MentorQA}) and the closest grounded or domain-specific care systems.}
\label{tab:positioning}
\end{table*}
